\chapter{Related work\label{cha:chapter3}}

\section{Satellite SIF Enhancement Techniques}
\label{sec:sif_enhancement}

Satellite observations of solar induced fluorescence (SIF) give information on vegetation photosynthesis activity but instrument characteristics in the form of information retrieval noise and uncertainties and satellite characteristics in the form of spatial footprint resolution and revisit times restrict their direct use for downstream tasks like GPP, crop yield experiments. Satellites with sensors with small footprints like OCO-2 and OCO-3 sample narrow and discontinuous ground tracks (each track is recorded on one date). Whereas wide swath instruments aboard satellites like TROPOMI \cite{kohler2018global} and GOME-2 \cite{anema2026long} provide more complete coverage but at a coarser spatial resolution. Therefore SIF enhancement methods have been developed to address these limitations by finding relationships between native SIF and spatially complete environmental variables, relationships between SIF observations or by improving retrieval from measured radiance. SIF enhancement techniques can be broadly grouped into three categories depending on their method of enhancement. \emph{Reconstruction and gap filling} methods estimate SIF at unobserved locations or dates and can extend a limited SIF record using an earlier predictor archive. \emph{Spatial downscaling} methods estimate how coarse observed or reconstructed SIF is distributed among finer subpixels. A third hybrid group combines both methods and therefore contains retrieval enhancements, forecasting (prediction) and geostatistical methods for enhancement. Identifying which method was used is important when interpreting the spatial resolution of the enhanced SIF product. For example, a value reported on a 500m grid is not necessarily a SIF observation at 500m because it was inferred from learned relationship assumptions \& auxiliary variables and not necessarily validated with real SIF samples at 500m resolution. Many of the methods reaggregate the fine pixel SIF estimates to the supervisory grid for validation. Therefore comparison of SIF enhancement methods considers how each method produces SIF and how it is evaluated.

\subsection{Reconstruction and Gap Filling}

Early globally contiguous SIF enhanced products confirmed that the data driven reconstruction strategies were suitable for the task. For instance, GOSIF \cite{li2019global} aggregates OCO-2 SIF at $0.05^{\circ}$ and learns its relationship with MODIS EVI, PAR, vapour pressure deficit (VPD), and air temperature using Cubist model trees. Using the spatially and temporally complete predictor grids, they were able to complete the OCO-2 gaps at $0.05^{\circ}$ and extend the record to year 2000. It showed that sparse OCO-2 SIF observations can support a global reconstruction. CSIF \cite{zhang2018global} also used the MODIS satellite data specifically four reflectance bands to reconstruct instantaneous daily SIF from raw OCO-2 SIF using a simple feed forward neural network. One important finding was that differences between CSIF and observed OCO-2 SIF could reveal short term environmental stress that the predictor reflectance alone couldn't capture. Instead of applying a single model across all biomes, Yu et al. SIFoco2\_005 product \cite{yu2019high}  fitted separate neural networks (ANNs) for each biome and 16 day time step, using seven-band MODIS reflectance to fill OCO-2 orbital gaps. The study found that subgrouping into separate biomes improved the distribution of training samples and the time specific models allowed the SIF reflectance relationship to vary with plant phenology and  stress. Not taking care of this biome and time stratification caused seasonal under-prediction and over-prediction and reduced drought sensitivity of the model. SIFoco2\_005 product was externally tested with air plane recorded SIF values (using the Chlorophyll Fluorescence Imaging Spectrometer (CFIS)) and thus provided a much stronger validation than same sensor testing. Though like the previous studies the resulting output was only viable at $0.05^{\circ}$, much coarser resolution than what would be needed for crop pure field scale studies. The spatially continuous TanSat product by Ma et al. \cite{ma2020generation} used a random forest model, MODIS optical variables and cosine of solar zenith angle for SIF gap filling. This study showed that cosine of solar zenith angle was a strong proxy for the amount and geometry of incoming sunlight (illumination), since SIF intensity depends partly on sunlight absorbed by vegetation this variable was important for SIF enhancement. 

The previous OCO-2 reconstructions mainly focused on spatial completion of the spatially sparse OCO-2 records, but reconstruction can also extend a temporally short satellite record over a longer historical period. TROPOMI provides near-global daily coverage of SIF at $0.05^{\circ}$, but its observations only became available in 2018 and still contain gaps caused by clouds and quality filtering, therefore temporal completion could be applied for this dataset. Chen et al. \cite{chen2022long} addresses this with their RTSIF product which used an XGBoost model to train on clear sky TROPOMI SIF with predictors; MODIS reflectance, land-surface temperature, land cover, PAR, C3/C4 vegetation fractions after which the trained model was applied to historical predictor dataset to extend the range of TROPOMI SIF historically. The variables chosen for modelling are important in the light use efficiency model \cite{liu2020estimating}, where reflectance characterizes canopy state and absorbed radiation, while temperature and vegetation type help represent changes in fluorescence efficiency. Compared to the previous studies, DOSIF pursues a finer temporal objective by producing daily, globally contiguous OCO-3-like SIF \cite{yu2025dosif}. It used a moving spatial-temporal window sampling strategy where a separate model was constructed for each day of the year using observations from nearby dates, multiple years, and geographically distinct sub-biomes. CatBoost model was found to be the best model for DOSIF's strategy. Again like \cite{yu2019high}, this sampling strategy outperformed a single global model indicating that relationships between SIF and reflectance vary across regions and stages of vegetation development (seasonal growth of vegetation). 

ST-LGBM \cite{shen2022spatiotemporal} like the previous OCO-2 SIF supervisory source studies tried to address the demanding problem of SIF gap filling large spatial gaps between OCO-2 track records. They used NIRv, PAR, temperature, vapour-pressure deficit, and land cover data in tabular format and therefore couldn't naturally contain neighbouring spatial information unlike in an image format. They addressed this by making spatial SIF factors from geographically nearby observations with similar vegetation state and temporal factors from observations at similar stages of the vegetation growth cycle in other years. This allowed the LightGBM model to use information from neighbouring regions (swaths) rather than relying only on variables observed at the same location and time without neighbourhood information. When they tested the LightGBM model without the feature engineered spatial and temporal factors on the OCO-2 SIF incomplete regions it performed much worse than the model with the spatial and temporal factors. This demonstrated that random validation may overestimate performance in areas without OCO-2 observation unless provided with the spatial and temporal factors. An even finer temporal reconstruction can also be done as demonstrated by He et al. \cite{he2025tracking}. Their study fit separate XGBoost models with OCO-3 SIF near 10:00 and 14:00 local time and GOCI reflectance bands as predictor variables producing continuous 500m resolution SIF maps for the morning and afternoon satellite overpasses. They normalized SIF using PAR and NIRv to produce fluorescence yield to reduce the influence of incoming sunlight and canopy structure and used the morning-to-afternoon normalized SIF differences as a drought indicator. The study showed that water and heat stress can be recorded in sub-daily observation (reduction in photosynthetic activity occurring within the same day, plants can close their stomata early during a heat stress), such an effect cannot be realised in a single SIF day observation. These reconstruction studies showed that SIF reconstruction methods have developed from simple models applied globally to models that consider location, time, nearby SIF observations and changes within a day. They also show that model performance depends on how the validation was done and whether the input variables capture short term changes in plant activity rather than only capturing physical properties of the vegetation canopy (optical variables like NDVI, NIRv etc.).  

\subsection{Spatial Downscaling}

Spatial downscaling means that a relationship learned between a coarse target variable (SIF) and a set of predictors can be evaluated using the same predictors at a high resolution to estimate the target variable at a finer resolution. Zhang et al. \cite{zhang2024spatial} used GOSIF \cite{li2019global} as the target trained with a random forest with predictors; NDVI, daytime land surface temperature at $0.05^{\circ}$ resolution after which the trained model was applied on 1km predictors to produce enhanced GOSIF at 1~km resolution. The enhanced GOSIF supports a stronger association with anomaly drought index and MODIS GPP than the original GOSIF. Similarly a related study instead uses a CNN model with NDVI, EVI, day and night time temperature to downscale GOSIF to approximately 1~km \cite{zhang2020downscaling}. Unlike the tabular data representation, the convolutional representation incorporates adjacent pixel context which was found to be useful for model fitting. Though it should be noted model evaluation mainly relied on reaggregated (to match their spatial resolution) enhanced GOSIF with original GOSIF and relationships with GPP, drought, or yield. These results as with almost all study results demonstrate only an application value and does not independently validate the enhanced SIF at the fine scale, as the best instrument recorded SIF resolution is around 1km$^2$ - 2km$^2$ (OCO-2, resolution varies with latitude). Also since GOSIF itself is an enhanced SIF product, its errors will carry over to any further enhancement experiments. Duveiller et al. downscale GOME-2 SIF using finer resolution NIRv, NDWI, and daytime land surface temperature \cite{duveiller2020spatially}. External validation with OCO-2 is used to choose the best model configuration. It was also successfully validated with TROPOMI after a small pixel level correction. Unlike methods mentioned earlier that try to predict SIF independently, this approach distributes the measured GOME-2 SIF within each coarse pixel to a finer resolution. 

Hu et al. \cite{hu2024spatially} BCSIF product first used a random forest model to predict high resolution TROPOMI SIF from spectral indices, illumination and environmental variables after which they estimated the difference between observed coarse resolution TROPOMI and predicted aggregated SIF and then redistributed this bias to a $0.005^{\circ}$ grid. This bias correction helps preserve the physiological information present in TROPOMI base SIF but not captured by the relationship between SIF and its predictors. Jin et al. \cite{jin2024downscaling} used random forest to predict high resolution SIF from coarse GOME SIF and predictors from AVHRR and ERA5. Residuals are then produced from the aggregated high resolution predictions versus the coarse GOME SIF. They then use ordinary kriging to estimate residuals in locations without observations based on residuals of nearby regions, after which the estimated residual is then added to the earlier high resolution random forest prediction thereby correcting its local errors. Their use of predictors from AVHRR and ERA5 helped SIF reconstruction before the MODIS era. SIFnet \cite{gensheimer2022convolutional} first coarsened TROPOMI SIF to $0.5^{\circ}$ and then trained a CNN (configured to learn the residuals) to recover the original $0.05^{\circ}$ SIF map from the coarsened SIF ($0.5^{\circ}$) and $0.05^{\circ}$ resolution secondary predictor variables (upscaled from original $0.005^{\circ}$), thereby learning a tenfold scale transformation. They assumed that this scale transformation can then be applied to $0.05^{\circ}$ native TROPOMI and $0.005^{\circ}$ secondary predictors to produce a 500m resolution SIF map. They validated the 500m resolution downscaled SIF by aggregating to OCO-2 and OCO-3 footprints, reporting RMSE = 0.21, R$^2$ = 0.57 and RMSE = 0.19, R$^2$ = 0.62 respectively. The parent coarsened TROPOMI SIF was the important feature followed by NIRv, Cosine of the solar zenith. 

TroDSIF \cite{chen2025improved} like the previous tabular methods uses random forest to predict SIF at 500m resolution, but instead of treating it as a final SIF map, the predictions are then used to determine how the original coarse TROPOMI SIF should be distributed among the smaller pixels by means of Gaussian redistribution. Therefore, the high resolution SIF map preserves the original TROPOMI measurement while adding spatial detail assuming that there was strong linear proportionality between the original SIF and predicted SIF (only then redistribution is valid). CNSIF combines long term SIF reconstruction with CNN based downscaling \cite{du2025cnsif}. The CNN uses GOSIF, NIRv, EVI, near-infrared reflectance, and land-surface temperature after which transfer learning is used to produce 500m estimates from Landsat data. It was validated on held out 2019 GOSIF samples rather than OCO-2 or OCO-3 SIF samples. Like the previous GOSIF enhanced studies, CNSIF will also inherit the uncertainties from GOSIF SIF estimates. HSIF \cite{zhang2023generating} instead of using a statistical learning method like the previous studies uses a physical based light use efficiency and radiative transfer method to determine how coarse TROPOMI SIF should be distributed to 1~km resolution. The relationships are determined using physical variables and equations, such as absorbed radiation, fluorescence efficiency, and canopy escape probability, rather than by fitting a machine learning model to labelled SIF samples. Overall the downscaling approaches show that high resolution predictor variables can add useful spatial detail to coarse SIF observations. But the reliability of this detail depends on whether the downscaling method preserves the original SIF measurement and if the downscaled SIF product was validated using high resolution original SIF measurements like from OCO-2 or OCO-3. 

\subsection{Other and Hybrid Enhancement Approaches}

Some SIF enhancement studies enhance SIF using hybrid approaches that use a mix of various techniques like reconstruction, gap filling, downscaling and geostatistical approaches. Pais et al. \cite{pais2025forecasting} firstly gap fills in the sparse OCO-2 SIF dataset with a random forest model using MODIS spectral indices, reflectance bands and land surface temperature variables as predictors, after which it then forecasts monthly and seasonal SIF using two earlier SIF maps, EVI, spatial and temporal variables. They compared various statistical and deep learning models, but an important point is errors from the first gap filling step can affect later forecasts. Li et al. \cite{li2025more} instead tries to denoise the original SIF retrievals by averaging matched OCO-2/3 observations within TROPOMI footprints and using them to train ANNs on TROPOMI light measurements at different wavelengths. This approach reduces noise in individual SIF retrievals and can therefore provide better training data for later reconstruction or downscaling. Tadić et al. \cite{tadic2024hybrid} used the previous CSIF enhanced SIF dataset to get an initial SIF estimate at every location, after which they then compared CSIF values with nearby OCO-2 SIF measurements. They then used kriging (geostatistical approach) to transfer differences between the CSIF and OCO-2 values to nearby local CSIF areas that have no OCO-2 measurements, thereby using the difference to improve the estimate. coSIF \cite{jacobson2023spatial} uses the cokriging method to model monthly OCO-2 SIF with atmospheric carbon dioxide from the following month. This method essentially learns how SIF values are related across the spatial setting and also how SIF is related to later changes in carbon dioxide. With this learned relationship, it can then use nearby SIF and carbon dioxide data to estimate missing SIF values. 

In conclusion, there is no single SIF enhancement strategy that can resolve all the limitations of satellite acquired SIF data. Global reconstructions of SIF can provide continuous SIF enhanced records but may smooth biome specific conditions. Studies have shown making separate models for each biome can improve the SIF enhancement performance. Non observation preservation downscaling approaches likewise can also produce high resolution SIF estimates but most of these methods use validation between these aggregated enhanced SIF maps with the coarse SIF value and strong relationship with GPP. These validation techniques cannot guarantee if the high resolution SIF map is accurate at the spatial support level because incorrect high resolution values can average to the correct coarse value, ie; the spatial distribution of the fine pixels in the coarse SIF can be wrong. Unlike the previous method, observation preserving and physical approaches can stay true to the original coarse SIF signal (ie; they average to original value) but again spatial aggregation of the fine pixels can still be inaccurate. Geostatistical methods (studies that used kriging method on OCO-2 SIF data) take care of local observation structure and SIF retrieval uncertainty. Though one critical issue with these methods is that their SIF gap-filling performance weakens as gaps between original SIF samples grow and nearby observations disappear. Using the OCO-2 dataset which is highly sparse spatially and temporally can greatly decrease  the viability of these methods on a lot of regions. All of these methods have their advantages and disadvantages, the usage depends on the availability of the respective predictor data, quality of the original SIF signals in the area of study.


\begingroup
\tiny
\setlength{\tabcolsep}{1pt}
\renewcommand{\arraystretch}{1.10}

\begin{longtable}{@{}
>{\raggedright\arraybackslash}p{0.110\linewidth}
>{\raggedright\arraybackslash}p{0.085\linewidth}
>{\raggedright\arraybackslash}p{0.180\linewidth}
>{\raggedright\arraybackslash}p{0.130\linewidth}
>{\raggedright\arraybackslash}p{0.120\linewidth}
>{\raggedright\arraybackslash}p{0.060\linewidth}
>{\raggedright\arraybackslash}p{0.145\linewidth}
>{\raggedright\arraybackslash}p{0.130\linewidth}@{}}

\caption{Summary of SIF enhancement methods for reconstruction, spatial downscaling, forecasting, retrieval improvement, and geostatistical gap filling. Many of the models are externally tested using GPP estimations. CFIS records airborne SIF while tower SIF is recorded by ground based towers. The $\dagger$ indicates an unofficial acronym.}
\label{tab:sif_enhancement_literature}\\
\hline
\textbf{Product} &
\textbf{SIF source} &
\textbf{Predictors} &
\textbf{Model family} &
\textbf{Spatio-Temporal resolution} &
\textbf{Period} &
\textbf{Internal validation metrics} &
\textbf{External validation metrics} \\
\hline
\endfirsthead

\multicolumn{8}{l}{\tablename\ \thetable{} - continued}\\
\hline
\textbf{Product} &
\textbf{SIF source} &
\textbf{Predictors} &
\textbf{Model family} &
\textbf{Spatial / temporal resolution} &
\textbf{Period} &
\textbf{Internal validation metrics} &
\textbf{External validation metrics and data} \\
\hline
\endhead

\hline
\multicolumn{8}{r}{Continued on next page}\\
\endfoot

\hline
\endlastfoot

\multicolumn{8}{@{}l}{\textbf{Reconstruction (no aggregation constraint)}} 
\vspace{2mm}\\

GOSIF \cite{li2019global} &
OCO-2 &
MODIS EVI; PAR, VPD, air temperature &
Cubist model trees &
Global; $0.05^{\circ}$; 8-day &
2000-2017 &
$R^2=0.79$; RMSE $=0.07$ &
FLUXNET GPP proxy: $R^2=0.73$ \\

TanSIF \cite{ma2020generation} &
TanSat &
MODIS reflectance, NDVI, $\cos(\mathrm{SZA})$, air temperature &
Random forest &
Global; $0.05^{\circ}$; 4-day &
2017-2019 &
$R^2=0.73$-$0.81$; RMSE $=0.30$-$0.32$ &
TROPOMI SIF: $R^2=0.73$ \\

DOSIF \cite{yu2025dosif} &
OCO-3 &
MODIS 7-band reflectance; biome and temporal context &
MSTWS + CatBoost &
Global; $0.05^{\circ}$; daily &
2001-2024 &
$R^2=0.81$-$0.85$; RMSE $=0.07$-$0.08$ &
Airborne CFIS: $R^2=0.75$; RMSE $=0.11$ \\

RTSIF \cite{chen2022long} &
TROPOMI &
MODIS reflectance, LST and land cover; PAR; C3/C4 fractions &
XGBoost &
Global; $0.05^{\circ}$; 8-day &
2001-2020 &
$R^2=0.907$; RMSE $=0.062$ &
PhotoSpec tower SIF: $R^2=0.754$-$0.84$ \\

ROSIF \cite{he2025tracking} &
OCO-3 &
Eight GOCI reflectance bands &
XGBoost &
Liaoning; 500 m; 10:00/14:00 &
2019-2020 &
$R^2=0.598$-$0.684$; RMSE $=0.348$-$0.388$ &
Drought proxies: SPEI $r=0.71$; SMZ $r=0.53$ \\

ST-LGBM \cite{shen2022spatiotemporal} &
OCO-2 &
NIRv, PAR, VPD, temperature, land cover; spatial/temporal SIF factors &
Constrained LightGBM &
Global; $0.05^{\circ}$; 8-day &
2014-2019 &
$R^2=0.79$-$0.83$; RMSE $=0.07$ &
TROPOMI SIF: annual $R^2=0.91$ \\

SIFoco2\_005 \cite{yu2019high} &
OCO-2 &
MODIS 7-band reflectance; biome/time stratification &
Feed-forward ANN &
Global; $0.05^{\circ}$; 16-day &
2014-2018 &
$R^2=0.83$; RMSE $=0.065$ &
Airborne CFIS outside swaths: $R^2=0.72$; slope $=0.96$ \\

CSIF \cite{zhang2018global} &
OCO-2 &
MODIS red, NIR, blue and green reflectance &
Feed-forward NN &
Global; $0.05^{\circ}$; 4-day &
2000-2017 &
$R^2=0.786$; RMSE $=0.177$ &
FLUXNET GPP proxy: median $R^2=0.64$; RMSE $=1.67$ \\

\hline
\multicolumn{8}{@{}l}{\textbf{Spatial downscaling}} \vspace{2mm}\\

RF-GOSIF\textsuperscript{$\dagger$} \cite{zhang2024spatial} &
GOSIF &
MODIS NDVI and daytime LST &
Random forest &
Henan; 1 km; monthly &
2001-2020 &
Reaggregated GOSIF: $r=0.74$ &
MODIS GPP proxy: $r=0.74$; crop yield: $r=0.89$-$0.93$ \\

G2DSIF\textsuperscript{$\dagger$} \cite{duveiller2020spatially} &
GOME-2 &
NIRv, NDWI and afternoon LST &
Semi-empirical LUE redistribution &
Global; $0.05^{\circ}$; 8-day &
2007-2018 &
- &
OCO-2 SIF: $r=0.813$-$0.825$; abs. bias $=0.0615$-$0.0630$ \\

BCSIF \cite{hu2024spatially} &
TROPOMI &
MODIS reflectance, NDWI, $\cos(\mathrm{SZA})$; temperature, VPD, soil moisture &
RF + bias correction &
China; $0.005^{\circ}$; 16-day &
2019-2020 &
$R^2=0.887$; RMSE $=0.0716$ &
ChinaSpec tower SIF: $R^2=0.798$; RMSE $=0.141$ \\

CNN-GOSIF\textsuperscript{$\dagger$} \cite{zhang2020downscaling} &
GOSIF &
NDVI, EVI, daytime and nighttime LST &
2D CNN &
Henan; $0.008^{\circ}$; monthly &
2013-2017 &
Reaggregated GOSIF: $r\leq0.9174$ &
MODIS GPP proxy: $r\leq0.8346$; Maize yield: $r=-0.84$ \\

DSIFRFK \cite{jin2024downscaling} &
GOME &
AVHRR NIRv, NDVI and NIR; temperature, radiation, precipitation &
RF + residual kriging &
East Asia; $0.05^{\circ}$; monthly &
1995-2003 &
$R^2=0.83$; RMSE $=0.08$ &
Other SIF products: $r=0.73$-$0.89$; tower GPP: $R^2=0.61$-$0.99$ \\

SIFnet \cite{gensheimer2022convolutional} &
TROPOMI &
Coarse SIF; MODIS reflectance/VIs; climate, soil, land cover, elevation &
Residual CNN &
CONUS; $0.005^{\circ}$; 16-day &
2018-2021 &
$R^2=0.92$; RMSE $=0.17$ &
OCO-2/3 footprints: $R^2:0.57$-$0.62$; RMSE $:0.19$-$0.21$ \\

TroDSIF \cite{chen2025improved} &
TROPOMI &
MODIS reflectance, NDVI, $\cos(\mathrm{SZA})$, air temperature &
RF + Gaussian redistribution &
Global; 500 m; 16-day &
2018-2021 &
Validation $R^2=0.8935$; RMSE $=0.064$; reaggregation $R^2=0.934$-$0.948$ &
Six ChinaSpec sites: RMSE $=0.104$-$0.223$; MAE $=0.077$-$0.163$ \\

HSIF \cite{zhang2023generating} &
TROPOMI &
Chlorophyll fPAR, fluorescence efficiency, escape probability, PAR, land cover &
Physical LUE/radiative-transfer model &
Global; 1 km; source-day/8-day &
2018-2020 &
- &
OCO-2 total SIF: $R^2=0.78$; RMSE $=1.33$; tower GPP proxy: $R^2=0.70$ \\

CNSIF \cite{du2025cnsif} &
GOSIF &
GOSIF, NIRv, EVI, NIR and LST from Landsat &
2D CNN + transfer learning &
China; 500 m; monthly &
2003-2022 &
$R^2=0.8566$; RMSE $=0.089$ &
Nine ChinaSpec sites: pooled $R^2=0.581$; RMSE $=0.152$; site $R^2=0.324$-$0.947$ \\

\hline
\multicolumn{8}{@{}l}{\textbf{Other and hybrid enhancement approaches}}\vspace{2mm} \\

OCO2-FSIF\textsuperscript{$\dagger$} \cite{pais2025forecasting} &
OCO-2 &
Two lagged SIF fields, EVI, spatial and temporal variables &
RF gap fill; Lasso/LightGBM forecast &
India; ${\sim}1$ km; monthly/seasonal &
2017-2019 &
RMSE $=0.0355$-$0.0594$; MAPE $=16.91$-$19.22\%$ &
- \\

ANN-TroSIF\textsuperscript{$\dagger$} \cite{li2025more} &
Matched OCO-2/3 &
TROPOMI radiance spectra at 743-758 nm &
Column-specific feed-forward ANNs &
Global; native TROPOMI footprint; daily &
2018-2024 &
$R^2=0.8535$; RMSE $=0.217$ &
Flux-tower GPP proxy: $R^2=0.45$-$0.79$ \\

KED-SIF\textsuperscript{$\dagger$} \cite{tadic2024hybrid} &
OCO-2 &
CSIF external drift and covariance of nearby OCO-2 SIF &
Kriging with external drift &
Global; $0.05^{\circ}$; 4-day &
2019 &
$R^2=0.8523$; RMSE $=0.1556$ &
- \\

coSIF \cite{jacobson2023spatial} &
OCO-2 &
SIF/XCO$_2$ covariance and spatial basis functions &
Multivariate cokriging &
North America; $0.05^{\circ}$; monthly &
2021 &
Spatial blocks: RASPE $=0.56$-$0.58$; bias $=-0.06$-$0.01$ &
- \\

\end{longtable}
\endgroup



\section{Remote Sensing-Based Crop Yield and GPP Estimation}
\label{sec:crop_yield_gpp}

A crop's yield can be dependent on various factors ranging from photosynthetic carbon uptake, respiration ratio, biomass allocation, effects from sudden shifts in weather conditions, soil and biological \& heat stress, all of which are factors during its growing season. Remote sensing from satellites can only capture some of these factors from optical radiometers and spectrometers. The signals that they can capture they do so usually in high temporal and spatial coverage. The following reviewed studies can be grouped into four approaches: SIF based yield estimation, optical and spectral index, multi source yield estimation and finally SIF based GPP estimation. Some of these groups can overlap in terms of scope as they first estimate the GPP and then derive the yield from its direct relationship as shown in Table \ref{tab:crop_yield_gpp_literature} while other studies use both SIF and spectral index predictors for estimation. The main SIF based yield estimation studies directly related reported yield (target) with seasonal (month wise) SIF using machine learning models, thereby showing data driven statistical approaches were suitable for estimating crop yields \cite{peng2020assessing} \cite{zheng2025modeling} \cite{joshi2023winter}. Peng et al. \cite{peng2020assessing} compared the yield predictive power of various SIF products: gap filled OCO-2, TROPOMI and much coarser GOME-2 SIF along with secondary variables like MODIS indices, land surface temperature for maize and soybean in the U.S Midwest. They showed that the higher resolution gap filled OCO-2 and TROPOMI SIF data outperformed GOME-2 by a significant margin and thus showed that the higher resolution SIF signals contain more crop pure SIF signals. They also noted NIRv (spectral index) can sometimes match SIF in performance depending on the crop type, validation method and training record length, though SIF still remained the dominant predictor variable for majority of the tests. Their tests included select regions in the North China Plain and CONUS (US), where a SIF based XGBoost reported a $R^2=0.87$ for the former while the CONUS study found that adding SIF to NIRv explained yield variation from 64\% to 69\%. 

Studies such as \cite{qiu2022monitoring}, \cite{song2018satellite} found that SIF is particularly informative when stress induced by drought, heatwave affects the photosynthesis process well before any visible canopy structure changes. As a reminder SIF is a faint light signal emitted by the plant's chlorophyll when absorbing sunlight during photosynthesis, while spectral indices like NIRv, NDVI, EVI measure plant canopy structure, greenness. The stress can affect the plant photosynthesis process before any changes can be seen in canopy. \cite{qiu2022monitoring} found that during the 2012 U.S Midwest drought yield declined by 25\%, where predictor variables GOSIF (enhanced SIF product) declined by 22\% ($R^2=0.91$) while NDVI, EVI and NIRv declined by only 4-10\%. Likewise \cite{song2018satellite} found that SIF detected reduced winter wheat growth activity in March well before NDVI and EVI showed any changes for the 2010 Indian heatwave. They also showed that SIF had a bigger response to yield than fPAR (both the variables are connected by the light use efficiency model \cite{liu2020estimating}), thereby indicating that the plant reduced its photosynthetic efficiency rather than only absorbing less sunlight. Though it should be noted that early stress detection does not guarantee accurate final yield estimation as SIF predicted a 13.9\% yield loss compared with the reported yield loss of 6\%. 

Unlike the previous studies, yield estimation studies like \cite{guan2016improving} \cite{he2020ground} \cite{kira2024scalable} \cite{xue2025lightweight} \cite{wang2025practical} first estimate Gross Primary Product (GPP) from SIF and then estimate crop yield from GPP after some corrections (also called process guided models). The common predictors used for these models include corrected canopy escape probability, C$_3$ and C$_4$ composition, respiration, NPP:GPP ratio, harvest index and harvested area. He et al. \cite{he2020ground} found that seasonal TROPOMI SIF successfully validated with ground based flux tower (eddy covariance tower) GPP measurements with a $R^2=0.89$, with the secondary predictor C$_3$/C$_4$ composition increasing productivity relationship from $R^2=0.72$ to 0.86. Kira et al. \cite{kira2024scalable} showed that a mechanistic light reaction model (non statistical based) was comparable to machine learning based techniques like ANN, random forest for U.S. maize yields. For regions like the Indo-Gangetic plain that had limited high quality remote sensing data, the mechanistic model in fact outperformed the machine learning models showing alternative ways of estimating GPP and crop yields. Another way of estimating yield was using $q_L$ (proportion of PSII reaction centres that are open. The centres help convert light into chemical energy. A higher $q_L$ thereby means more capacity for photosynthesis) and SIF to estimate plant photosynthesis. Wang et al. \cite{wang2025practical} used this method to explain about 78\% of differences in county level maize and soybean yields $R^2 = 0.78$. These mechanistic models are more interpretable than regression models as they use the direct physical process for yield estimation but they still depend on crop specific parameters and assumptions that are used to convert photosynthetic uptake to yield. Therefore the choice between direct SIF-yield methods and SIF-GPP-yield methods depends on data availability. The machine learning models can learn complex relationships when there are long, reliable yield records whereas the process guided models need fewer yield records (but require more plant specific data like respiration, crop allocation, $q_L$).

Older spectral index only approaches like \cite{vannoppen2021estimating} \cite{chandel2019yield} \cite{hunt2019high} \cite{vallentin2022suitability} tried to estimate crop yield from canopy greenness and structure, chlorophyll, plant water content (each variable has a related spectral index). Seasonal NDVI summaries were found to be weak winter wheat yield predictors compared to precipitation data observed during wheat tillering and flowering stage ($R^2=0.66$) for Belgium winter wheat yields \cite{vannoppen2021estimating}. In contrast \cite{chandel2019yield} found that NDVI and NDWI (recorded during heading stage; the stage when the grain head first emerges) explained roughly 95-96\% of grain yield variation on a small set of farms in central India, but it was not validated on regional scale. Hunt et al. \cite{hunt2019high} used Sentinel-2 spectral indices and showed a yield prediction performance of RMSE from 0.66 to 0.61~t~ha$^{-1}$ on U.K. wheat yield data. Vallentin et al. \cite{vallentin2022suitability} conducted a yield estimation study for Northeast Germany (13 years) and showed that spectral indices correlated with yield ranging from almost 0 to 0.94, with a median of only 0.19, thereby showing an overall weak correlation between spectral indices and yield. They also found that the correct data acquisition timing, field heterogeneity (high resolution spectral indices at 10m can capture spatial details in farms that were heterogeneous) mattered more for high spectral index-yield correlation than choice of a particular spectral index. Overall the spectral index only approaches showed a relatively weaker yield estimation performance than the main SIF based approaches. 

Multi source studies combine spectral index data (which describe plant canopy characteristics), weather and soil data (growth constraints) and crop-growth simulations (which provide data for water stress) to estimate yields as shown in studies like \cite{goihl2024crop} \cite{marszalek2022prediction} \cite{brandt2024ensemble} \cite{zhao2020predicting} \cite{srivastava2022winter}. In southern Germany, raw Sentinel-2 bands along with precipitation and evapotranspiration data explained 84\% of winter wheat yield variance \cite{marszalek2022prediction}. Likewise \cite{zhao2020predicting} showed for a northeastern Australian set, Sentinel-2 index predictors along with crop-growth simulation water stress index and chlorophyll indices could explain 93\% of wheat yield differences. Brandt et al. used six multi source estimators to conduct yield estimation at the parcel (crop field) level with roughly 140,000-155,000 parcel samples available per year until 2019. They reported a $R^2=0.74$ at parcel level and 0.79-0.89 at the district level (NUTS 3) after aggregation. Yield performance from these studies cannot be directly compared because their target resolution ranges from parcel (field) level to county or district averages. Some of the studies tested on a random test sample which could contain same year and region samples and therefore cannot test genuine extrapolation from training to test. Studies that tested on aggregated county or district samples can report a much higher $R^2$ as aggregation reduces local noise. 

The remaining yield studies like \cite{hu2020detecting}, \cite{ma2025gpp}, \cite{shekhar2022well} focused only on estimating GPP and not the final crop yield from SIF, but it still provides a physiological link between SIF and biomass / yield production. Hu et al. \cite{hu2020detecting} found that using local relationships to down scale GOME-2 SIF represented regional and seasonal changes in GPP better than using one global relationship to downscale GOME-2 SIF. Ma et al. \cite{ma2025gpp} used a transfer learning approach to combine the strengths of two datasets: SIF-derived GPP and Eddy-covariance (ground based FLUX tower) GPP. Their transfer learning model was first trained on the SIF derived GPP data to learn general relationships between GPP and predictors like NIRv, leaf area index, incoming solar radiation, temperature and vapour pressure deficit after which the pretrained model was then fine tuned using GPP measurements from 164 FLUX tower sites. This gave an improvement of $R^2$ = 0.773 compared with the pretrained model showing  $R^2$ = 0.641. Though such a study is limited in geographic scope and is highly local, transferability to other regions depends on availability of nearby FLUX tower data. It should also be noted that an accurate GPP estimation does not necessarily determine the crop yield. Other factors like respiration, timing of carbon uptake, allocation to non grain biomass, harvest index affect the final conversion of GPP into yield. These studies and the earlier main SIF based approaches confirm that SIF is the important variable for crop yield / GPP estimation as it can directly measure photosynthetic activity of the plant.  It is also useful for detecting heat stress in crops and a good proxy for identifying drought conditions. Other auxiliary data like spectral indices, weather, soil, crop composition can add useful information with SIF to better explain the yield process. Common limitations of these methods are low resolution satellite records (mainly for SIF data) giving poor local / parcel level information, temporal and spatial gaps in the data, uncertain target crop area and yield statistics which can affect the modelling process, and spatial aggregation of yield statistics causing accuracy gains / noise reduction compared to local samples. These findings therefore support that spatially enhanced SIF along with high resolution spectral and environmental variables are useful for crop yield or GPP estimations. 



\section{Comparative Analysis and Research Gap}
\label{sec:comparative_analysis_research_gap}

This section briefly identifies important findings and methodological gaps in existing SIF enhancement methods that can be addressed by future studies. Spatially and temporally continuous SIF was shown that it can be reconstructed from sparse satellite observations using reflectance, vegetation indices, radiation, temperature, and other environmental variables using either statistical data driven approaches or mechanistic models depending on the availability of data. Most of the discussed methods that preserve the original coarse SIF observation as part of their learning process provide more consistent results than completely unconstrained prediction, though this can be challenging with sparse OCO-2 and OCO-3 datasets. Discrepancy between output enhanced SIF resolution and available SIF products at high resolution for validation is a challenge that needs to be addressed. Some studies have used airborne \cite{yu2019high}, \cite{yu2025dosif}  or FLUX tower based SIF recordings \cite{hu2024spatially}, \cite{chen2025improved}, \cite{du2025cnsif} (high resolution) to validate their enhanced SIF product, but the validation remains viable only at the local level and cannot be generalized on a much larger scale. Across all these methods, the most important variable was the parent SIF observation (usually coarse) when one exists followed by secondary predictor variables like optical variables like reflectance, NIRv, radiation variables like PAR, fPAR, land surface temperature, vapour pressure deficit and soil moisture content. Parent SIF was only available in continuous format for TROPOMI and GOSIF \cite{li2019global} and some of the kriging methods could use spatially sparse OCO-2 samples, but this was only reliable when their regional cluster density was high. Land cover classification and crop composition which were available at high spatial resolution was also found to be important for decomposing the coarser SIF signal to a finer scale. Majority of the methods used a tabular model as reliable large sample training for the convolutional models can be challenging. Both tabular and convolutional models were able to reliably enhance coarse SIF measurements, with tabular methods often including neighbouring information through feature engineered spatial factors. The CNN methods \cite{gensheimer2022convolutional}, \cite{du2025cnsif}, \cite{zhang2020downscaling} could naturally learn neighbourhood patterns, field boundaries and additional spatial context that cannot be reliably represented in tabular format. However a CNN model cannot recover physiological information that is absent from its predictors. Therefore predictor quality, spatial support and external validation design mattered more than choosing a convolutional or tabular model. No single study tested if model architecture matters when provided with the same data, spatial and temporal scope, therefore an interesting point of research could be to address this.

Using OCO-2 as the base target support (unlike TROPOMI, GOME-2) can be challenging especially in gap filling contexts because of its narrow recording track, spatially wide swaths between each track. For a single year OCO-2 SIF data, there will be large regions without parent OCO-2 SIF observations to draw any context of general magnitude of SIF in this region. Though this can be mitigated by combining multiple year data at the same location, but in turn lose year specific characteristics. One major research gap is that majority of the methods output enhanced SIF is at \(0.05^\circ\) resolution with a few methods reporting 500m, no SIF product exists at a finer spatial resolution than 500m. For crop yield studies in regions with small farm sizes (usually under 60 hectares), a higher than 500m resolution SIF product is necessary to acquire crop pure SIF signals. Most of the models' highest resolution predictors were the MODIS reflectance bands at 250m-500m resolution. Sentinel-2 provides reflectance band data in 10m, 20m, 60m resolution across its 13 bands. As per the literature review, OCO-2 reconstruction with very high resolution predictors like Sentinel-2 predictors has not been attempted and therefore remains a promising research direction.  It also remains insufficiently tested whether convolutional models outperform tabular models when both use the same Sentinel-2 information. As previously discussed, SIF offers greater crop yield predictor power than other secondary data, therefore it remains to be tested if isolating crop pure SIF from mixed SIF footprints can support more accurate crop yield estimations \cite{qiu2022monitoring,song2018satellite,peng2020assessing}. 

\clearpage
% crop yield table
\newlength{\cropyieldtablewidth}
\setlength{\cropyieldtablewidth}{\linewidth}

\begingroup
\tiny
\setlength{\tabcolsep}{1pt}
\renewcommand{\arraystretch}{1.10}
\setlength{\LTcapwidth}{\cropyieldtablewidth}

\begin{longtable}{@{}
>{\raggedright\arraybackslash}p{0.110\cropyieldtablewidth}
>{\raggedright\arraybackslash}p{0.100\cropyieldtablewidth}
>{\raggedright\arraybackslash}p{0.150\cropyieldtablewidth}
>{\raggedright\arraybackslash}p{0.140\cropyieldtablewidth}
>{\raggedright\arraybackslash}p{0.120\cropyieldtablewidth}
>{\raggedright\arraybackslash}p{0.140\cropyieldtablewidth}
>{\raggedright\arraybackslash}p{0.060\cropyieldtablewidth}
>{\raggedright\arraybackslash}p{0.120\cropyieldtablewidth}@{}}

\caption{Summary of remote-sensing-based crop yield, crop productivity, and GPP studies. All RMSE values are reported in t ha$^{-1}$ units}.
\label{tab:crop_yield_gpp_literature}\\
\hline
\textbf{Reference} &
\textbf{Prediction target} &
\textbf{Remote-sensing inputs} &
\textbf{Additional predictors} &
\textbf{Model} &
\textbf{Study area} &
\textbf{Period} &
\textbf{Performance} \\
\hline
\endfirsthead

\multicolumn{8}{l}{\tablename\ \thetable{} - continued}\\
\hline
\textbf{Reference} &
\textbf{Prediction target} &
\textbf{Remote-sensing inputs} &
\textbf{Additional predictors} &
\textbf{Model / approach} &
\textbf{Study area / scale} &
\textbf{Period} &
\textbf{Reported performance} \\
\hline
\endhead

\hline
\multicolumn{8}{r}{Continued on next page}\\
\endfoot

\hline
\endlastfoot

\multicolumn{8}{@{}l}{\textbf{SIF-based yield, productivity, and stress estimation}}
\vspace{2mm}\\

Peng et al. \cite{peng2020assessing} &
Maize / Soybean &
GOME-2/downscaled SIF; NDVI; MODIS GPP &
County effects &
Fixed-effects regression &
United States; county &
2015--2018 &
$R^2=0.85 - 0.89$; RMSE $=0.85$ \\

Guan et al. \cite{guan2016improving} &
Maize / Soybean &
GOME-2 SIF &
C3/C4 type; respiration; temperature &
Process-guided conversion &
United States; county &
2007-2012 &
$R^2=0.43 - 0.72$ \\

Zheng et al. \cite{zheng2025modeling} &
Winter-wheat yield &
SIF; NIRv, NDVI, EVI, EVI2 &
Climate and soil &
XGBoost; RF; ANN; SVM &
North China Plain; municipality &
2007-2020 &
$R^2=0.87$; RMSE $=0.352$\\

Xue et al. \cite{xue2025lightweight} &
Wheat production &
Satellite SIF &
Temperature; soil water; VPD &
Process-guided $q_L$ model &
Australia; state &
2019-2022 &
$R^2=0.86$-$0.91$\\

Kira et al. \cite{kira2024scalable} &
Maize / wheat &
OCO-2 SIF; MODIS NDVI, NIR, fPAR &
PAR; NPP:GPP; area; harvest index &
MLR; ANN; RF &
U.S. counties; Indian districts &
2015-2020 &
$R^2=0.53$-$0.69$ \\

He et al. \cite{he2020ground} &
Maize / Soybean  &
PhotoSpec and TROPOMI SIF &
GPP; planted area; C3/C4 fractions &
Cross-scale regression &
U.S. Maize Belt; tower/county &
- &
$R^2=0.72$-$0.89$ \\

Wang et al. \cite{wang2025practical} &
Maize / Soybean yield &
Satellite SIF &
Temperature; VPD; soil moisture; phenology &
Process-guided $q_L$ model &
U.S. Midwest; county &
2018-2023 &
$R^2=0.76$-$0.81$ \\

Qiu et al. \cite{qiu2022monitoring} &
Maize / Soybean &
GOSIF; GOME-2 SIF; NDVI, EVI, NIRv &
Tower GPP; yield statistics &
Anomaly regression &
U.S. Midwest; regional &
2008-2018 &
$R^2=0.91$ for GOSIF \\

Joshi et al. \cite{joshi2023winter} &
Winter-wheat yield &
SIF and three VIs &
Climate and soil &
XGBoost; RF &
CONUS; county &
14 years &
$R^2=0.69$ \\

Song et al. \cite{song2018satellite} &
Heat stress &
SIF, SIF yield, NDVI, EVI &
Temperature; absorbed light; yield &
Anomaly comparison &
Indo-Gangetic Plains &
2010 &
Yield loss: observed $6\%$; SIF $13.9\%$ \\

\hline
\multicolumn{8}{@{}l}{\textbf{Optical and spectral-index yield estimation}}
\vspace{2mm}\\

Vannoppen et al. \cite{vannoppen2021estimating} &
Winter-wheat yield &
Maximum NDVI &
Monthly precipitation &
Random forest &
Belgium; field &
2016-2018 &
$R^2=0.66$  \\

Hunt et al. \cite{hunt2019high} &
Wheat yield &
Sentinel-2 &
Weather; terrain; soil moisture &
Random forest &
United Kingdom; 10 m/field &
- &
RMSE $= ~0.63$ \\

Vallentin et al. \cite{vallentin2022suitability} &
Cereal &
Six sensors; 15 spectral indices &
Timing; soil; terrain; field properties &
Correlation analysis &
Northeast Germany; field &
13 years &
$r\approx0$-$0.94$; median $0.19$ \\

Chandel et al. \cite{chandel2019yield} &
Wheat yield &
NDVI, NDWI, SAVI, NDNI &
Nitrogen treatment; growth stage &
Stage-specific regression &
Central India &
2014-2017 &
$R^2=0.95$-$0.96$ \\

\hline
\multicolumn{8}{@{}l}{\textbf{Multi-source and temporal yield models}}
\vspace{2mm}\\

Goihl et al. \cite{goihl2024crop} &
Six crop yields &
- &
Weather; historical yield &
Random forest &
Saxony &
2015-2022 &
Wheat RRMSE $=2.88$-$20.56\%$ \\

Marszalek et al. \cite{marszalek2022prediction} &
Winter-wheat yield &
Sentinel-2 bands; NDRE, REIP, NDWI &
Rainfall; evapotranspiration &
Linear regression; RF &
Southern Germany; field &
2016-2018 &
$R^2=0.84$; RMSE $=0.56$ \\

Brandt et al. \cite{brandt2024ensemble} &
Three crop yields &
Multi-source EO variables &
Weather; soil; parcels &
Six-model ensembles &
Two German states; parcel/district &
2019-2022 &
$R^2=0.66$-$0.89$; nRMSE $=7.2$-$16.9\%$ \\

Zhao et al. \cite{zhao2020predicting} &
Wheat yield &
Sentinel-2 OSAVI, NDVI, CI, NDRE &
Crop-model water stress &
Combined regression &
Northeast Australia; field &
2016-2017 &
Test $R^2=0.93$; RMSE $=0.64$ \\

Srivastava et al. \cite{srivastava2022winter} &
Winter-wheat yield &
- &
Weekly weather; soil; phenology &
1D CNN &
Germany; county &
1999-2019 &
RMSE reduction $=7$-$14\%$ \\



\hline
\multicolumn{8}{@{}l}{\textbf{SIF-based GPP estimation and evaluation}}
\vspace{2mm}\\

Hu et al. \cite{hu2020detecting} &
Regional GPP &
GOME-2 SIF; MODIS NDVI, fPAR, LST &
Soil-moisture index &
Moving-window regression &
Flux-tower regions &
- &
- \\

Ma et al. \cite{ma2025gpp} &
Global GPP &
Three SIF-derived GPP sources &
Eddy-covariance GPP &
Transfer learning &
Global; tower network &
- &
$R^2$: $0.036$-$0.132$ \\

Shekhar et al. \cite{shekhar2022well} &
GPP evaluation &
CSIF, GOSIF, LUE-SIF, HSIF; VIs &
FLUXNET/ICOS GPP &
Statistical comparison &
Global/Europe; tower &
2007-2018 &
SIF generally exceeded NDVI/EVI \\

\end{longtable}
\endgroup
